\chapter{PHƯƠNG PHÁP}
\label{chap:methodology}

\section{Tổng quan}

% Hình~\ref{fig:architecture} minh họa kiến trúc tổng thể của hệ thống phân tích cảm xúc được đề xuất.
%
% \begin{figure}[H]
% \centering
% \includegraphics[width=0.9\textwidth]{figures/architecture.png}
% \caption{Kiến trúc tổng thể của hệ thống phân tích cảm xúc được đề xuất.}
% \label{fig:architecture}
% \end{figure}
%
% TODO: Thêm hình architecture.png vào thư mục figures/ và bỏ comment các dòng trên

\section{Xử lý dữ liệu}

\subsection{Chuẩn hóa văn bản}

Chuẩn hóa văn bản bao gồm:
\begin{itemize}
    \item Chuyển đổi thành chữ thường
    \item Xóa các ký tự đặc biệt và URL
    \item Mở rộng viết tắt
    \item Sửa lỗi chính tả
\end{itemize}

\subsection{Tách từ}

Sử dụng phương pháp tách từ phù hợp cho văn bản tiếng Việt để tách câu thành các từ riêng lẻ.

\section{Trích xuất đặc trưng}

\subsection{Word Embeddings}

Sử dụng pre-trained word embeddings để biểu diễn văn bản dưới dạng các vector số.

\section{Kiến trúc mô hình}

\subsection{Mô tả mô hình}

Mô hình của chúng tôi bao gồm vài lớp:

\begin{enumerate}
    \item Input layer: Nhận văn bản đã xử lý
    \item Embedding layer: Chuyển tokens thành vector dense
    \item Encoding layer: Xử lý chuỗi
    \item Classification layer: Xuất ra xác suất cảm xúc
\end{enumerate}

\subsection{Công thức toán học}

Cho một chuỗi đầu vào $X = \{x_1, x_2, \ldots, x_n\}$, mô hình tính toán:

\begin{equation}
h = f(X; \theta)
\end{equation}

trong đó $f$ là hàm encoding và $\theta$ đại diện cho các tham số mô hình.

Dự đoán cảm xúc được tính toán:

\begin{equation}
\hat{y} = \arg\max_i \text{softmax}(W h + b)_i
\end{equation}

\section{Huấn luyện}

\subsection{Hàm mất mát}

Sử dụng cross-entropy loss:

\begin{equation}
\mathcal{L} = -\sum_{i=1}^{C} y_i \log(\hat{y}_i)
\end{equation}

trong đó $C$ là số lượng lớp cảm xúc.

\subsection{Tối ưu hóa}

Mô hình được tối ưu hóa bằng Adam optimizer với learning rate $\alpha = 0.001$.
